% =============================================================
%  01_modello_mentale.tex
%  Capitolo 1 — Il modello mentale di Lisp
% =============================================================

\chapter{Il modello mentale di Lisp}

\begin{intuizione}
Lisp non è un linguaggio da \emph{imparare}, è un linguaggio da
\emph{capire}. La sintassi si impara in dieci minuti. Il modello mentale
richiede qualche ora. Questo capitolo serve a costruirlo.

\medskip
Vieni da C, Java, Python: sai già scrivere cicli e assegnamenti.
Qui imparerai qualcosa di diverso - il modo funzionale di pensare
ai problemi. È quello che ti darà il vantaggio reale dallo studio di Lisp.
L'imperativo lo conosci già; il funzionale è la parte nuova.
\end{intuizione}

\vspace{4pt}

% ================================================================
\section{Tutto è un'espressione}
% ================================================================

In C, Python o Java distingui nettamente \emph{istruzioni} ed
\emph{espressioni}: \texttt{if} è un'istruzione, \texttt{a + b} è
un'espressione. In Lisp questa distinzione non esiste.

\begin{intuizione}
\textbf{In Lisp tutto è un'espressione e ogni espressione produce un valore.}\\
Non ci sono istruzioni. Non ci sono statement. Solo espressioni.
\end{intuizione}

\argomento{S-expression}

La struttura fondamentale del linguaggio si chiama \textbf{S-expression}
(symbolic expression). Ha due forme:

\begin{sintassi}
atomo          ; un valore singolo: numero, stringa, simbolo\\
(f a1 a2 ...)  ; una lista: primo elemento = operatore/funzione
\end{sintassi}

\mnota{S-expression (symbolic expression) — introdotta da McCarthy nel 1960.}

Esempi di atomi:

\begin{esempio}[Atomi]
\begin{lstlisting}
42           ; numero intero
3.14         ; numero in virgola mobile
"ciao"       ; stringa
t            ; vero (true)
nil          ; falso / lista vuota
pi           ; costante predefinita (~3.14159)
\end{lstlisting}
\end{esempio}

Esempi di liste (chiamata di funzione):

\begin{esempio}[Liste come chiamate]
\begin{lstlisting}
(+ 3 5)        ; => 8
(* 4 2)        ; => 8
(+ 1 2 3 4)    ; => 10   (+ accetta qualsiasi numero di argomenti)
(max 7 3 9 2)  ; => 9
\end{lstlisting}
\end{esempio}

\begin{attenzione}
L'operatore va \textbf{sempre prima} degli argomenti. Non esiste
\texttt{3 + 5}: si scrive \texttt{(+ 3 5)}.\\[4pt]
Se vieni da Python o C, questa è l'unica cosa che richiede
adattamento. Dopo pochi minuti diventa automatica.
\end{attenzione}

\seprule

\argomento{Annidamento}

Le espressioni si annidano liberamente. Il risultato di una
espressione è l'argomento di un'altra:

\begin{esempio}[Espressioni annidate]
\begin{lstlisting}
(+ (* 3 4) (* 2 5))   ; => 22
; equivale a: 3*4 + 2*5 = 12 + 10

(sqrt (+ (* 3 3) (* 4 4)))   ; => 5.0
; equivale a: sqrt(3^2 + 4^2)
\end{lstlisting}
\end{esempio}

\mnota{Nessuna precedenza tra operatori: le parentesi definiscono tutto. Zero ambiguita'.}

\begin{nota}
Non esiste precedenza tra operatori: le parentesi definiscono esplicitamente
l'ordine di valutazione. Questo è una \emph{feature}, non un limite:
il codice non è mai ambiguo.
\end{nota}

% ================================================================
\section{Il REPL: il tuo laboratorio}
% ================================================================

\begin{concetto}{REPL — Read Eval Print Loop}
Il REPL è l'ambiente interattivo di Lisp. Il flusso è:

\medskip
\begin{center}
\textbf{Read} $\to$ legge un'espressione\\[3pt]
\textbf{Eval} $\to$ la valuta\\[3pt]
\textbf{Print} $\to$ stampa il risultato\\[3pt]
\textbf{Loop} $\to$ ricomincia
\end{center}

\medskip
Non scrivi un programma, lo compili e lo esegui. Scrivi
un'espressione, premi Invio, vedi il risultato. Subito.
\end{concetto}

\argomento{Workflow nel REPL}

La sessione tipica:

\begin{esempio}[Sessione REPL]
\begin{lstlisting}
; Avvia Allegro CL Free Express e digita:
CL-USER> (+ 1 2)
3                          ; il REPL stampa il risultato

CL-USER> (* 6 7)
42

CL-USER> (defun quadrato (x) (* x x))
QUADRATO                   ; Lisp conferma la definizione

CL-USER> (quadrato 5)
25

CL-USER> (quadrato (+ 2 1))
9
\end{lstlisting}
\end{esempio}

\begin{workbook}
Apri il REPL adesso e testa queste espressioni prima di andare avanti.
Prova a variare i numeri, annida le espressioni, rompi qualcosa e leggi
l'errore.
\end{workbook}

\argomento{Caricare un file}

Per lavori più lunghi scrivi il codice in un file \texttt{.lisp}
e lo carichi nel REPL:

\begin{sintassi}
:ld nomefile.lisp    ; Allegro CL (shorthand)\\
(load "nomefile.lisp")  ; ANSI standard, funziona ovunque
\end{sintassi}

\mnota{In Allegro: \fn{:ld}, \fn{:cl} (compile+load), \fn{:trace}, \fn{:untrace}.}

\begin{nota}
Flusso di lavoro consigliato: scrivi 3-5 righe nel file, carica,
testa nel REPL, correggi. Non aspettare di avere 100 righe prima
di testare.
\end{nota}

% ================================================================
\section{Valutazione: le regole del gioco}
% ================================================================

\begin{concetto}{Come Lisp valuta un'espressione}
Regole di valutazione — semplici e uniformi:

\medskip
\begin{enumerate}
  \item \textbf{Atomo auto-valutante}: numeri, stringhe, \texttt{t}, \texttt{nil}
        restituiscono se stessi.
  \item \textbf{Simbolo}: viene cercato nell'ambiente corrente e restituisce
        il suo valore.
  \item \textbf{Lista \texttt{(f a1 a2 ...)}}: prima si valutano tutti gli
        argomenti (da sinistra a destra), poi si chiama \texttt{f} con i risultati.
  \item \textbf{Forma speciale}: alcune parole chiave (\texttt{if}, \texttt{let},
        \texttt{defun}, \texttt{quote}\ldots) hanno regole di valutazione proprie.
\end{enumerate}
\end{concetto}

\argomento{Quote: bloccare la valutazione}

A volte vuoi che un simbolo o una lista \emph{non} venga valutata
ma trattata come dato. Si usa \texttt{quote} (o l'abbreviazione \texttt{'}):

\begin{esempio}[Quote]
\begin{lstlisting}
(+ 1 2)      ; => 3  (valuta: chiama +)
'(+ 1 2)     ; => (+ 1 2)  (non valuta: restituisce la lista)

'ciao        ; => CIAO  (simbolo come dato, non come variabile)
(quote x)    ; equivalente a 'x
\end{lstlisting}
\end{esempio}

\begin{intuizione}
\textbf{Codice e dati hanno la stessa forma.}
Una lista \texttt{(+ 1 2)} è sia un'espressione da valutare
sia un dato (una lista di tre elementi). Il \texttt{quote} sceglie
il secondo caso. Questa simmetria — chiamata \emph{homoiconicity} —
è il fondamento delle macro.
\end{intuizione}

% ================================================================
\section{Differenza con il paradigma imperativo}
% ================================================================

Vieni da C, Java, Python. La differenza più profonda non è la sintassi:
è come pensi al programma.

\begin{confronto}
\textbf{In C o Java} un programma è una sequenza di \emph{istruzioni}
che modificano uno stato.
Si \emph{fa} qualcosa: si assegna, si itera, si modifica.

\medskip
\textbf{In Common Lisp} un programma è una composizione di
\emph{espressioni} che producono valori.
Non ci sono istruzioni: ogni forma restituisce qualcosa.
\fn{if} restituisce un valore. \fn{let} restituisce un valore.
Anche una funzione che stampa a schermo restituisce un valore.
\end{confronto}

\begin{esempio}[Stesso calcolo: imperativo vs Lisp]
\begin{lstlisting}
;;; Python
x = 3
y = x * 2
z = y + 1   ; z vale 7

;;; Common Lisp: tutto e' un'espressione
(let* ((x 3)
       (y (* x 2))
       (z (+ y 1)))
  z)           ; => 7

;;; O direttamente, senza variabili:
(+ (* 3 2) 1)  ; => 7
\end{lstlisting}
\end{esempio}

La tabella seguente riassume i punti di rottura principali:

\begin{confronto}
\renewcommand{\arraystretch}{1.4}
\small
\begin{tabular}{@{}p{5.8cm}p{6.0cm}@{}}
\rowcolor{boxCfFrame!12}
\textbf{\color{boxCfFrame!80!black}Imperativo (C / Java)} &
\textbf{\color{csBlue}Common Lisp} \\[2pt]
%
\texttt{result = f(x)} \newline
{\footnotesize\color{black!55}assegnazione} &
\texttt{(let ((result (f x))) \ldots)} \newline
{\footnotesize\color{black!55}binding locale} \\
\hline
\texttt{if (...) \{\} else \{\}} \newline
{\footnotesize\color{black!55}istruzione, non restituisce valore} &
\texttt{(if cond vero falso)} \newline
{\footnotesize\color{black!55}espressione con valore} \\
\hline
\texttt{for (int i=0; i<n; i++)} \newline
{\footnotesize\color{black!55}loop imperativo} &
\texttt{(loop for i from 0 below n)} \newline
{\footnotesize\color{black!55}o ricorsione} \\
\hline
\texttt{void f(x) \{ ... \}} \newline
{\footnotesize\color{black!55}procedura senza return} &
{\footnotesize\color{black!55}Non esiste: ogni funzione} \newline
{\footnotesize\color{black!55}restituisce sempre un valore} \\
\hline
\texttt{f = \&myfunc} \newline
{\footnotesize\color{black!55}puntatore a funzione (raro)} &
\texttt{\#'f} \ o \ \texttt{(lambda (x) \ldots)} \newline
{\footnotesize\color{black!55}funzioni come valori normali} \\
\end{tabular}
\end{confronto}

\seprule

\begin{intuizione}
\textbf{Common Lisp è multi-paradigma}: funzionale, imperativo, a oggetti (CLOS).
Non è Haskell: \fn{setf}, array mutabili, hash table e side-effect
sono strumenti normali, da usare quando ha senso.

\keyline{}{Il funzionale è lo stile naturale, non l'unico consentito.}
\end{intuizione}

\begin{attenzione}
\textbf{Parti dal funzionale, non dall'assegnazione.}\\
\fn{setf} è l'operatore di assegnazione: esiste ed è uno strumento legittimo.
Ma se lo usi come default fin dall'inizio, stai replicando
i pattern di C senza guadagnare nulla dal cambio di linguaggio.
\begin{lstlisting}
;; Imperativo (da evitare all'inizio)
(setf x (+ x 1))

;; Funzionale: calcola un nuovo valore
(let ((x-nuovo (+ x 1)))
  ...)
\end{lstlisting}
\end{attenzione}

% ================================================================
\section{Definire funzioni: \texttt{defun}}
% ================================================================

\begin{concetto}{\texttt{defun}: definire una funzione}
\begin{sintassi}
(defun nome (param1 param2 ...) corpo)
\end{sintassi}
\begin{itemize}
  \item \texttt{nome}: il simbolo con cui richiamare la funzione
  \item \texttt{(param1 ...)}: lista dei parametri
  \item \texttt{corpo}: una o più espressioni; la funzione restituisce
        il valore dell'\textbf{ultima}
\end{itemize}
\end{concetto}

\begin{esempio}[Prime funzioni]
\begin{lstlisting}
(defun add (a b)
  (+ a b))

(add 3 5)    ; => 8
(add 1.5 2)  ; => 3.5


(defun quadrato (x)
  (* x x))

(quadrato 4)   ; => 16
(quadrato -3)  ; => 9


;; Corpo con piu' espressioni: conta solo l'ultima
(defun saluta (nome)
  (format t "Ciao, ~a!~%" nome)   ; effetto collaterale: stampa
  nil)                             ; valore restituito: nil

(saluta "Alice")
; stampa: Ciao, Alice!
; => NIL
\end{lstlisting}
\end{esempio}

\begin{nota}
\texttt{format t "~a~\%"} stampa il valore di \texttt{~a} seguito da
newline (\texttt{~\%}). La \texttt{t} indica lo standard output.
Non preoccuparti ora dei dettagli di \texttt{format}: verrà trattato
nel capitolo sull'I/O.
\end{nota}

\argomento{Ricorsione e iterazione}

La ricorsione è particolarmente naturale per attraversare strutture
ricorsive come liste e alberi: il codice rispecchia direttamente
la struttura del dato.
Per iterazioni imperative, Common Lisp offre costrutti potenti
come \fn{loop}, \fn{dotimes} e \fn{dolist}.
In codice professionale si usano entrambi, scegliendo
il più leggibile caso per caso.

\begin{esempio}[Prima funzione ricorsiva]
\begin{lstlisting}
(defun fattoriale (n)
  (if (= n 0)
      1                         ; caso base
      (* n (fattoriale (- n 1)))))  ; caso ricorsivo

(fattoriale 5)   ; => 120
(fattoriale 0)   ; => 1
\end{lstlisting}
\end{esempio}

\mnota{Ogni funzione ricorsiva ha un caso base (stop) e uno ricorsivo (avanzamento).}

\begin{intuizione}
\textbf{Come leggere una funzione ricorsiva:}
Non seguire l'esecuzione passo per passo nella testa — ti perdi.
Invece: \emph{"se n è 0, restituisce 1; altrimenti, n per il
fattoriale del resto"}.
Fidati della definizione. Il computer gestisce lo stack.
\end{intuizione}

\begin{workbook}
\textbf{Esercizi 1--4} (Fondamenti: Somma, Geometria, Conversioni, Ipotenusa)\\[4pt]
Sono le prime funzioni: pura sintassi \texttt{defun} e aritmetica.
Scrivile nel REPL, non in un file. Obiettivo: prendere confidenza con
la notazione prefissa e con la struttura \texttt{(defun nome (params) corpo)}.
\end{workbook}

% ================================================================
\section{Nomi, simboli e case}
% ================================================================

\begin{concetto}{Simboli in Common Lisp}
Un \textbf{simbolo} è un nome: può essere una variabile, una funzione,
una costante. Common Lisp converte tutto in maiuscolo internamente,
ma la convenzione di scrittura è \textbf{kebab-case} minuscolo:
\end{concetto}

\begin{esempio}[Convenzioni di naming]
\begin{lstlisting}
;; Nomi di funzione: kebab-case
(defun calcola-area (r) (* pi r r))
(defun converti-celsius->fahrenheit (c) (+ (* c 9/5) 32))

;; Predicati (restituiscono vero/falso): terminano con -p
(defun pari-p (n) (= (mod n 2) 0))
(defun vuoto-p (lst) (null lst))

;; Variabili globali: asterischi (per convenzione)
;; defparameter definisce una variabile globale con un valore iniziale
(defparameter *pi-greco* 3.14159)
(defconstant +max-iter+ 1000)

;; Il case non conta per il lettore:
(defun Foo (x) x)   ; identico a (defun foo (x) x)
\end{lstlisting}
\end{esempio}

\begin{center}
\renewcommand{\arraystretch}{1.5}
\begin{tabular}{@{}lll@{}}
\toprule
\textbf{Tipo} & \textbf{Convenzione} & \textbf{Esempio} \\
\midrule
Funzione       & kebab-case          & \texttt{calcola-area} \\
Predicato      & kebab-case + \texttt{-p} & \texttt{pari-p}, \texttt{null} \\
Variabile globale & \texttt{*nome*}  & \texttt{*debug*} \\
Costante       & \texttt{+nome+}     & \texttt{+max-size+} \\
Parametro locale & kebab-case        & \texttt{lista-input} \\
\bottomrule
\end{tabular}
\end{center}

\begin{nota}
Le convenzioni non sono imposte dal compilatore, ma sono universalmente
rispettate nel codice Lisp professionale. Adottarle subito ti permette
di leggere codice altrui senza sforzo.
\end{nota}

% ================================================================
\section{Gestire gli errori nel REPL}
% ================================================================

Prima o poi (subito) il REPL mostra un errore. Non farti prendere dal
panico: il debugger di Lisp è potente e informativo.

\begin{esempio}[Errori tipici e come leggerli]
\begin{lstlisting}
;; Parentesi mancante: il REPL aspetta
CL-USER> (+ 1 2
               ; cursore lampeggia, si aspetta )
               )   ; digita la parentesi che manca
3

;; Funzione non definita
CL-USER> (foo 42)
; Error: The function FOO is undefined.
; Digita :reset per tornare al prompt

;; Tipo sbagliato
CL-USER> (+ 1 "ciao")
; Error: "ciao" is not of type NUMBER.

;; Troppi / troppo pochi argomenti
CL-USER> (add 1)
; Error: ADD requires exactly 2 arguments.
\end{lstlisting}
\end{esempio}

\begin{sintassi}
:reset    ; (Allegro) torna al prompt pulito dopo un errore\\
:bt       ; (Allegro) mostra lo stack trace\\
:q        ; esce dal debugger
\end{sintassi}

\begin{workbook}
\textbf{Esercizi 5--11} (Fondamenti: Distanza, Pari/dispari, Divisibilita', Segno,
Massimo, Triangolo, Anno bisestile)\\[4pt]
Questi esercizi richiedono \texttt{let}, \texttt{cond} e \texttt{mod}.
Prima di scriverli, leggi la sezione \emph{Condizioni e flusso}
(Capitolo~2) oppure consulta la sezione corrispondente nel workbook
Quick Reference.\\[4pt]
\textbf{Obiettivo:} consolidare \texttt{defun}, cominciare a usare
i condizionali, capire quando usare \texttt{let} per i valori intermedi.
\end{workbook}

% ================================================================
\section{Il ciclo di sviluppo professionale}
% ================================================================

\begin{concetto}{Come lavora un Lisper}
Non si scrive il programma intero e poi si compila. Si costruisce
\textbf{incrementalmente}, pezzo per pezzo, con il REPL come
amico fidato:

\medskip
\begin{enumerate}
  \item Scrivi una funzione piccola (5-10 righe max)
  \item Caricala nel REPL (\texttt{:ld file.lisp})
  \item Testala con 2-3 casi, incluso un caso limite
  \item Se funziona, vai avanti. Se no, correggi e ricarica.
  \item Componi le funzioni piccole in funzioni più grandi
\end{enumerate}

\medskip
Il REPL non è un accessorio: è \textbf{l'ambiente di sviluppo}.
\end{concetto}

\argomento{Trace: vedere l'esecuzione}

\texttt{:trace} è uno strumento prezioso per capire cosa fa una
funzione ricorsiva:

\begin{esempio}[Trace di una funzione ricorsiva]
\begin{lstlisting}
CL-USER> :trace fattoriale
; Tracing FATTORIALE

CL-USER> (fattoriale 4)
0: (FATTORIALE 4)
  1: (FATTORIALE 3)
    2: (FATTORIALE 2)
      3: (FATTORIALE 1)
        4: (FATTORIALE 0)
        4: FATTORIALE returned 1
      3: FATTORIALE returned 1
    2: FATTORIALE returned 2
  1: FATTORIALE returned 6
0: FATTORIALE returned 24
24

CL-USER> :untrace fattoriale   ; disattiva il trace
\end{lstlisting}
\end{esempio}

\mnota{\fn{:trace} mostra ogni chiamata ricorsiva con indentazione = profondita'.}

\begin{intuizione}
Usa \texttt{:trace} ogni volta che una funzione ricorsiva non si
comporta come ti aspetti. Vedere le chiamate annidate risolve il
90\% dei bug di ricorsione.
\end{intuizione}

% ================================================================
\section*{Riepilogo del capitolo}
% ================================================================

\begin{tcolorbox}[
  enhanced,
  colback=csBlue!5,
  colframe=csBlue!40,
  arc=4pt,
  boxrule=0.6pt,
  left=14pt, right=14pt, top=12pt, bottom=12pt,
  before skip=16pt,
  after skip=16pt
]
\textbf{Quello che hai imparato:}

\medskip
\begin{itemize}
  \item \textbf{S-expression}: la struttura uniforme di tutto il codice Lisp.
        Atom | \texttt{(operatore arg1 arg2 ...)}
  \item \textbf{Valutazione}: atomi si auto-valutano; liste chiamano funzioni;
        \texttt{quote} blocca la valutazione.
  \item \textbf{REPL}: il laboratorio interattivo. Scrivi, valuta, correggi.
  \item \textbf{defun}: come definire funzioni. Il corpo è un'espressione,
        non una sequenza di istruzioni.
  \item \textbf{Ricorsione}: naturale per strutture ricorsive (liste, alberi).
        Per iterazioni imperative: \fn{loop}, \fn{dotimes}, \fn{dolist}.
  \item \textbf{Convenzioni}: kebab-case, predicati con \texttt{-p},
        globali con \texttt{*asterischi*}.
  \item \textbf{Ciclo di sviluppo}: incrementale, con il REPL.
        Funzioni piccole, test immediati, composizione.
\end{itemize}

\medskip
\textbf{Cosa viene dopo:}\\
Capitolo~2, \emph{Funzioni e dati fondamentali}: tutti i tipi di dato,
le funzioni numeriche, le stringhe, \texttt{let}/\texttt{let*},
parametri opzionali e \texttt{\&rest}.

\medskip
\textbf{La direzione del percorso:}\\
I capitoli 1–4 ti chiedono di pensare in modo funzionale.
Resistenza garantita: il tuo istinto cercherà di scrivere cicli
e assegnamenti come sai fare. Tienilo a bada.
Dal capitolo~5 in poi vedrai anche la mutabilità usata dove serve,
come scelta consapevole - non come default.
\end{tcolorbox}

\begin{workbook}
\textbf{Prima di andare al Capitolo~2:}\\[4pt]
Completa gli esercizi \textbf{1--11} (capitolo \emph{Fondamenti}).\\[4pt]
Essenziali con quanto visto qui: \textbf{1, 6}
(pura \fn{defun} e aritmetica di base).\\[4pt]
Gli esercizi \textbf{4, 8} richiedono \fn{let*} e \fn{cond},
trattati nel Capitolo~2: puoi tentarli subito o tornare dopo
aver letto le prime sezioni di quel capitolo.\\[4pt]
Se hai difficolta' con un esercizio \textbf{$\bigstar\bigstar\bigstar$} o superiore,
consulta il suggerimento corrispondente nel workbook prima della soluzione.
\end{workbook}
